\documentclass[11pt,a4paper]{article}
\usepackage[T1]{fontenc}
\usepackage{lmodern,amsmath,amssymb,amsthm}
\usepackage[margin=23mm,top=19.5mm,bottom=22mm]{geometry}
\usepackage{microtype}
\usepackage{fancyvrb}
% Keep a rebuild independent of the host clock and absolute source path.
\pdfinfoomitdate=1
\pdftrailerid{}
\pdfsuppressptexinfo=15
\usepackage[hidelinks,pdflang={en},
  pdfauthor={Astra 6.2},
  pdftitle={\{aa,abab,baaa,b\} is a smallest code not included in any finite maximal code},
  pdfsubject={Formalized in Lean by Opus 5.5. Supervised by Christophe Cordero.
    Discovered on 26 November 2026.}]{hyperref}
\newtheorem{theorem}{Theorem}
\newtheorem{lemma}[theorem]{Lemma}
\newtheorem{proposition}[theorem]{Proposition}
\newcommand{\head}[1]{\par\noindent\textbf{#1}\enspace}
\newcommand{\eps}{\varepsilon}
\newcommand{\Z}{\mathbb{Z}}
\DeclareMathOperator{\tr}{tr}
\DeclareMathOperator{\adj}{adj}
\DeclareMathOperator{\rev}{rev}
\setlength{\parskip}{.2em}
\begin{document}
\begin{center}
  {\LARGE\bfseries\boldmath $\{aa,abab,baaa,b\}$ is a smallest code\\[2pt]not included in any
  finite maximal code\par}
\end{center}
\noindent\begin{tabular*}{\linewidth}{@{}l@{\extracolsep{\fill}}lll@{}}
{\small\sffamily Author} & {\small\sffamily Lean formalization} & {\small\sffamily Supervisor} &
  {\small\sffamily Date of discovery} \\
Astra 6.2 & Opus 5.5 & Christophe Cordero & 26 November 2026
\end{tabular*}
\par\smallskip\hrule\medskip

Let $A$ be a finite alphabet. Write $A^*$ for its words, $\eps$ for the empty word, and
$A^+=A^*\setminus\{\eps\}$. A set $Y\subseteq A^+$ is a \emph{code} if distinct finite
sequences of its elements have distinct concatenations. A \emph{finite maximal completion} of $Y$
is a finite code $M\supseteq Y$ maximal under inclusion among all codes on $A$. It is
\emph{positively factorizing} if there are finite $L,R\subseteq A^*$ such that every word has a
unique expression $\ell m_1\cdots m_jr$ with $j\ge0$, $\ell\in L$, $r\in R$ and $m_i\in M$. For
finite $Y$ put $\lambda(Y)=\max_{y\in Y}|y|$, with $\lambda(\emptyset)=0$, and
$\tau(Y)=\sum_{y\in Y}|y|$. On the empty alphabet the empty code is itself maximal and positively
factorizing with $L=R=\{\eps\}$.

\begin{theorem}[Simultaneous minimality]\label{thm:main}
Every finite code $Y$ on a finite alphabet has a positively factorizing finite maximal completion
on that alphabet whenever
\[
  |Y|\le3\qquad\text{or}\qquad\lambda(Y)\le3\qquad\text{or}\qquad\tau(Y)\le10.
\]
For any distinct letters $a,b$, the code $X=\{aa,abab,baaa,b\}$ has no finite maximal completion
on any finite alphabet containing $a,b$. It attains the three minima $4$, $4$ and $11$
simultaneously.
\end{theorem}

\section{A criterion on finitely many periodic words}

For a language $C$, let $C^*$ consist of all concatenations of its members, including $\eps$. A
word $w$ is \emph{framed} by $C$ if $lwr\in C^*$ for some $l,r\in C^*$. If $C\subseteq M$ and $M$
is a code, every framed word of $M$ belongs to $C$. Indeed, a parsing of $lwr$ over $C$ must
agree with the parsing over $M$ that uses $w$ as one word.

Fix $d\ge2$ and choose integers $j_r\ge0$ with $j_r\equiv r\pmod d$ for $0\le r<d$ and $j_0=0$.
Put
\[
  B=\{a^d\}\cup\{ba^{j_r}: 0\le r<d\},\qquad G=\{k\ge0: k<j_{k\bmod d}\}.
\]
In $B^*$ a word $ba^k$ must start with some $ba^{j_r}$ and continue in blocks of $a^d$. Thus $G$ is
precisely the finite set of gaps $k$ for which $ba^k\notin B^*$.

For $k\ge0$ let $v_k=a^kb$. On vertices $0,\dots,k$ put an $a$-edge of weight $1$ from $i$ to
$i+1$ for $i<k$, and a $b$-edge of weight $t$ from $k$ to $0$. Write $U,V$ for the two adjacency
matrices (rows are starting vertices) and $\rho_k$ for substitution $a\mapsto U$, $b\mapsto V$. For
a finite set $C$ let $\chi_C=\sum_{w\in C}w$ be its noncommutative polynomial.

\begin{theorem}[Finite-cycle criterion]\label{thm:criterion}
Let $Y\subseteq\{a,b\}^+$ be a finite code containing $B$ and assume $G\ne\emptyset$. Suppose
that for each $k\in G$, with $v=v_k$, the following conditions hold.
\begin{enumerate}
  \item[(i)] Every finite factor of the two-sided periodic word $v^{\Z}$ is framed by $Y$.
  \item[(ii)] For some $h\ge1$, $v^h\in Y^*$ and $\det(I-\rho_k(\chi_Y))=1-t^h$.
  \item[(iii)] For every $0<i<h$, $v^i\notin Y^*$ and there is $c_i\in Y^*$ with $c_iv^i\in Y^*$.
\end{enumerate}
Then $Y$ has no finite maximal completion on any finite alphabet containing $a,b$. These
conditions are decidable from the given data.
\end{theorem}

\begin{proof}
For a finite set $C\subseteq A^+$, Reutenauer's factorization theorem
\cite[\S2, corollary, p.~172]{Reutenauer} says that $C$ is a maximal code on $A$ if and only if
\[
  1-\chi_C=P(1-\chi_A)S,\qquad P,S\in\Z\langle A\rangle.
\]
The coefficients may have either sign. The constant coefficients have product $1$, so we may
normalize both to $1$. Suppose a finite maximal code contains $Y$. Apply this identity and send
every letter outside $\Sigma=\{a,b\}$ to zero. Writing $M$ for its restriction to $\Sigma^+$, we
have a code $M\supseteq Y$ and integral polynomials satisfying
\begin{equation}\label{eq:fact}
  1-\chi_M=P(1-a-b)S.
\end{equation}

\head{The gap polynomial.}If $w$ avoids every $ba^gb$ with $g\in G$, it is framed by $B$. Indeed,
choose a multiple $N$ of $d$ greater than every $j_r$. In $(ba^N)wa^N$ the initial and final runs
following a $b$ at the padded ends have length at least $N$, and every intervening $b$-to-$b$ gap
in $w$ lies outside $G$. Hence each block $ba^m$ belongs to $B^*$, as do both pads. This forces
every such word of $M$ into $B$. In particular the words of $M$ with at most one $b$ are exactly
those of $B$.

Write $P_0(x)=P(x,0)$ and $p(x,y)=\sum_{i,j\ge0}[a^iba^j]P\,x^iy^j$, where $[w]$ extracts the
coefficient of $w$. Use $S_0,s$ similarly and put $J_d(x)=1+x+\dots+x^{d-1}$. Comparing words
with zero and one occurrence of $b$ in \eqref{eq:fact} gives
\begin{equation}\label{eq:onebee}
\begin{aligned}
  P_0(x)S_0(x)&=J_d(x),\\
  \textstyle\sum_r y^{j_r}&=P_0(x)S_0(y)+(y-1)p(x,y)S_0(y)+(x-1)P_0(x)s(x,y).
\end{aligned}
\end{equation}
The three terms in the second line place the unique $b$ in the middle, left and right factors,
respectively. At $y=1$ the polynomial $P_0(x)$ divides the nonzero integer $d$, so its constant
term forces $P_0=1$. Now $S_0=J_d$. At $x=1$, \eqref{eq:onebee} yields
\begin{equation}\label{eq:gaps}
  p(1,y)=\frac{\sum_r(y^{j_r}-y^r)}{y^d-1}=\sum_{g\in G}y^g,\qquad p(1,1)=|G|.
\end{equation}
The last polynomial equality sums the missing exponents in each residue class. Since $b\in M$,
each $b^m$ with $m>1$ already parses as $m$ copies of $b$. Thus $b$ is the only pure $b$-word in
$M$. Substituting $a=0$ in \eqref{eq:fact} therefore gives $P(0,b)S(0,b)=1$, hence $P(0,b)=1$.

\head{The periodic determinants.}Write $D_k=\det\rho_k(P)$ and $H_k=\det\rho_k(S)$. They have
constant term $1$, since at $t=0$ their matrices are $I$ and $J_d(U)$ with $U$ strictly upper
triangular. Multilinearity of the determinant gives $[t]\det(I+tT+O(t^2))=\tr T$. A path reading
$a^iba^j$ is closed precisely when $i+j=k$, in which case it starts at $j$ and is unique. Thus
\begin{equation}\label{eq:trace}
  [t]D_k=\sum_{i+j=k}[a^iba^j]P.
\end{equation}

For the column vectors $e_i$ of the standard basis, $V=te_ke_0^{\mathsf T}$. Factoring out an
invertible $W$ and expanding in columns gives $\det(W-xy^{\mathsf T})=\det(W)(1-y^{\mathsf
T}W^{-1}x)$, since terms using two proportional columns vanish. Applied with $W=I-U$, this gives
$\det(I-U-V)=1-t$. Taking determinants in \eqref{eq:fact}, we obtain
\begin{equation}\label{eq:det}
  (1-t)D_kH_k=\det(I-\rho_k(\chi_M)).
\end{equation}
If $k\notin G$, every path avoids the forbidden gaps, so the readable words in $M$ and $B$
coincide. Their matrix is $U^d+V\sum_rU^{j_r}$. The same rank-one calculation, using
$(I-U^d)^{-1}=\sum_{q\ge0}U^{dq}$, gives
\[
  \det(I-\rho_k(\chi_M))=1-t\,\#\{(r,q):0\le r<d,\ q\ge0,\ j_r+dq=k\}=1-t.
\]
Thus $D_k=1$ for all $k\notin G$.

\head{The side of the factorization that carries a period.}Choose $K>\deg P$ and form the series
$F=\sum_{w\in\Sigma^*}1_{M^*}(wb^K)\,w$. For any language $T$ we use $\chi_T$ also for its
characteristic series. In the formal series ring, \eqref{eq:fact} inverts to
$\chi_{\Sigma^*}=S\chi_{M^*}P$ because $M$ is a code. For $|w|\ge\deg S$, the coefficient of
$wb^K$ can only use the empty suffix from $P$, by $P(0,b)=1$, and every contributing prefix from
$S$ lies in $w$. Consequently $SF=\chi_{\Sigma^*}+E$ for a polynomial $E$. The substitution
$\rho_k$ extends to these series because a path using $n$ letters $b$ has length at most
$(k+1)n+k$, so every coefficient is a finite sum. It sends $\chi_{\Sigma^*}$ to $(I-U-V)^{-1}$.
Multiplication by $\adj(\rho_k(S))$ and then by $1-t=\det(I-U-V)$ shows that
\begin{equation}\label{eq:poly}
  (1-t)H_k\,\rho_k(F)\quad\text{has polynomial entries.}
\end{equation}

Now fix $k\in G$. By (i), the words of $M$ readable on this cycle are exactly the readable words of
$Y$. Hence \eqref{eq:det} and (ii) give $(1-t)D_kH_k=1-t^h$. Moreover
\[
  v^nb^K\in M^*\quad\Longleftrightarrow\quad h\mid n.
\]
Multiples of $h$ follow from (ii) and $b\in Y$. Otherwise write $n=qh+i$, $0<i<h$. A $Y$-parsing
of $c_iv^i$ has no cut after $c_i$, since such a cut would give a $Y$-parsing of $v^i$. Extend it
by $v^{qh}b^K$. A parsing of $v^nb^K$ over $M$, preceded by a $Y$-parsing of $c_i$, would give a
cut at that very position, contradicting uniqueness in $M^*$. The closed paths at vertex $0$ read
exactly $v^n$, with weight $t^n$. Therefore $(\rho_k(F))_{00}=1/(1-t^h)$. By \eqref{eq:poly},
$1-t^h$ divides $(1-t)H_k$. The preceding determinant identity now makes $D_k$ a unit in
$\Z[t]$, and $D_k(0)=1$ gives $D_k=1$.

Summing \eqref{eq:trace} over $k\ge0$ thus yields $p(1,1)=0$. The sum on its right is finite since
$P$ is a polynomial. This contradicts \eqref{eq:gaps}.

\head{Effectiveness.}A finite automaton for $Y^*$ consists of a loop spelling each word at a common
initial and final vertex, with separate interior vertices. Determinize it by recording the set of
current vertices. To recognize framed words, allow as starts its reachable accepting states, and as
ends all states from which some word of $Y^*$ reaches an accepting state. These ends are computed
by a product with the automaton for $Y^*$. With all vertices as starts and ends, the $v$-cycle
recognizes its factors. Inclusion in the framed language is decided by searching its product with
the complement of the complete deterministic framed automaton for an accepting path. The
determinant determines $h$ if it is of the required form. The automaton for $Y^*$ tests the
required powers. A witness $c_i$ exists precisely when the transition on $v^i$ takes some
reachable accepting state to an accepting state.
\end{proof}

\section{The two-parameter family}

For finite sets of nonempty words put
\[
  \partial(C,T)=\{r\ne\eps: ur\in T\text{ for some }u\in C\}.
\]
Let $D(C)$ be the least set $D$ with $\partial(C,C)\cup\partial(D,C)\cup\partial(C,D)\subseteq D$.
All these residuals are proper nonempty suffixes of words of $C$, so the closure is finite.

\begin{lemma}[Residual test]\label{lem:residual}
A finite set $C\subseteq A^+$ is a code if and only if $D(C)\cap C=\emptyset$.
\end{lemma}

\begin{proof}
After cancelling common initial words of two distinct parsings of the same word, the two unequal
first words must be prefix-comparable. Their difference is an initial residual. At each subsequent
step extend the shorter parsing by its next word. The comparison either reaches equality when the
residual is that word, or gives one of the two strict quotients in the closure. Conversely each
residual is realized by two partial parsings with different first words, one concatenation
extending the other by the residual. This holds initially and is preserved by each comparison. A
residual in $C$ therefore closes an ambiguity.
\end{proof}

\begin{theorem}[A family with no finite maximal completion]\label{thm:family}
For all $d,s\ge2$, writing $u=ab$, the set
\[
  Y_{d,s}=\{a^d,ba^{d+1},u^s\}\cup\{ba^j:2\le j<d\}\cup\{bu^j:0\le j\le s-2\}
\]
is a code with no finite maximal completion on any finite alphabet containing $a,b$. The same
holds after reversing every word.
\end{theorem}

\begin{proof}
Choose $j_1=d+1$ and $j_r=r$ for $r\ne1$ in Theorem~\ref{thm:criterion}. Thus $G=\{1\}$. The
initial residuals lie among $a^e$ for $1\le e\le d+1$, $e\ne d$, and $u^i$ for $1\le i<s$. Indeed
the exponents of the words $ba^{j_r}$ have distinct residues modulo $d$, the words $bu^j$ are
nested, and $ba^j$ ($j\ge2$) differs from $bu^i$ ($i\ge1$) at the third letter. The two words
starting with $a$ differ at their second letter. Adjoining the single residual $bu^{s-1}$ gives a
set closed under all comparisons with a codeword, as the table shows.
\begin{center}
\begin{tabular}{lll}
Residual & Possible next residuals & Comparison word\\\hline
$a^e$ & $a^{|d-e|}$ & $a^d$\\
$a$ ($e=1$) & $bu^{s-1}$ & $u^s$\\
$u^i$ & $u^{s-i}$ & $u^s$\\
$bu^{s-1}$ & $u^{s-1-j}$ ($0\le j\le s-2$) & $bu^j$
\end{tabular}
\end{center}
The first two letters distinguish all other comparisons except those beginning with $ba$ in the
last row, where the third letter distinguishes the two forms. None of the listed residuals is a
codeword. Lemma~\ref{lem:residual} proves unique decipherability.

Set $c=ba^d$. For $n=qs+i$, $0\le i<s$, the two possible remainders give
\[
  cu^n=\begin{cases}(b)(a^d)(u^s)^q,& i=0,\\ (ba^{d+1})(bu^{i-1})(u^s)^q,& 1\le i<s.\end{cases}
\]
Thus $cu^n\in Y_{d,s}^*$ for every $n\ge0$, and $ca=ba^{d+1}\in Y_{d,s}$. For any nonempty
alternating word $w$ there are $\alpha\in\{\eps,a\}$ and $\beta\in\{\eps,b\}$ with $\alpha
w\beta=u^n$. The words $c\alpha$, $\beta$ frame $w$. For $0<i<s$ neither $a^d$ nor $u^s$ can begin
a parsing of $u^i$, and $cu^i\in Y_{d,s}^*$. On the $ab$-cycle the readable codewords are $u^s$
and $bu^j$ for $0\le j\le s-2$, giving
\[
  \rho_1(\chi_{Y_{d,s}})=\begin{pmatrix}t^s&0\\t+t^2+\dots+t^{s-1}&0\end{pmatrix}.
\]
Its required determinant is $1-t^s$, and $u^s\in Y_{d,s}$. All hypotheses of
Theorem~\ref{thm:criterion} hold with $h=s$. Reversal bijects parsings by reversing their order,
and preserves inclusions. It therefore preserves codes and maximality.
\end{proof}

\section{Positive completions for short words}

We first record the algebraic form of positive factorization. Let $P,S$ be finite polynomials on
$A$ with nonnegative integer coefficients and constant coefficient $1$. If
\begin{equation}\label{eq:cert}
  F=1+P(\chi_A-1)S
\end{equation}
has nonnegative coefficients, its support is a positively factorizing finite maximal code.
Indeed, $F$ has zero constant coefficient, and inversion gives
\[
  S\Big(\sum_{j\ge0}F^j\Big)P=(1-\chi_A)^{-1}=\chi_{A^*}.
\]
Each coefficient on the left counts choices with multiplicities. Taking the other factors empty
shows that $P$, $S$, $F$ all have coefficients at most one. The same identity gives the required
unique factorization. Maximality follows from Reutenauer's equivalence above. Conversely, if every
word has the unique expression $\ell m_1\cdots m_jr$ from the definition, then $\eps\in L\cap R$ by
factoring $\eps$, and inversion of the characteristic-series identity yields \eqref{eq:cert} with
$F=\chi_M$, $P=\chi_R$ and $S=\chi_L$. We call \eqref{eq:cert} a \emph{positive certificate} in
this situation.

\begin{proposition}\label{prop:short}
Every code $C\subseteq A\cup A^2\cup A^3$ has a positively factorizing finite maximal completion
with word lengths at most seven.
\end{proposition}

\begin{proof}
Write $D=D(C)$ and partition the letters into $A\cap C$, $V=A\cap D$ and $U=A\setminus(C\cup D)$.
Lemma~\ref{lem:residual} gives $C\cap D=\emptyset$. The closure rules imply, for $x,r\ne\eps$,
\begin{equation}\label{eq:rules}
  x\in C\cup D,\ xr\in C\ \Longrightarrow\ r\in D,\qquad x\in C,\ xr\in D\ \Longrightarrow\ r\in D.
\end{equation}
Use the sets
\[
  H=UA\cup AV,\qquad P_2=H\setminus(C\cup D),\qquad S_2=H\cap D.
\]
Complete the two factors by the sets of triples
\begin{align*}
  P_3&=\{xyz: xy\in P_2,\ x\notin U,\ z\in U,\ yz\notin D\},\\
  S_3&=\{xyz: yz\in S_2,\ z\notin V,\ x\in V,\ xy\in D\}.
\end{align*}
Take $P,S$ to be the characteristic polynomials of $\{\eps\}\cup U\cup P_2\cup P_3$ and
$\{\eps\}\cup V\cup S_2\cup S_3$. We show that \eqref{eq:cert} has nonnegative coefficients and
contains $C$ in its support. In the coefficient calculations write $[T]$ for the indicator of a
condition $T$.

\head{Words of lengths one to three.}The coefficient of $x$ is $1-[x\in U]-[x\in V]=[x\in C]$. For
$xy$ the coefficient is $[xy\in H]-[xy\in P_2]-[xy\in S_2]$, since
$[x\in U]+[y\in V]-[x\in U][y\in V]=[xy\in H]$. Every $xy\in C$ belongs to $H$, since $x\notin U$
forces $y\in V$ by \eqref{eq:rules}. Thus the coefficient is $[xy\in C]$.

Before the triples $P_3,S_3$ are included, the coefficient of $xyz$ is
\begin{equation}\label{eq:three}
  [z\notin V][xy\in P_2]+[x\notin U][yz\in S_2]+[x\in U][z\in V].
\end{equation}
For $xyz\in C$, the two implications
\[
  z\notin V\ \Longrightarrow\ xy\notin C\cup D,\qquad x\notin U\ \Longrightarrow\ yz\in D
\]
follow from \eqref{eq:rules}. If $x\in U$ or $z\in V$, the last summand in \eqref{eq:three} is $1$
when both hold. When just one holds, the first or second summand is $1$, respectively. If neither
holds, $y\notin C$ because $yz\in D$ and $z\notin D$. Thus $y\in V$ or $y\in U$, putting $xy$ in
$P_2$ or $yz$ in $S_2$, respectively. Adding $P_3,S_3$ subtracts one on each set. They are disjoint
since $xy\in P_2$ on $P_3$ and $xy\in D$ on $S_3$. The two implications exclude $C$ from both sets.
On $P_3$ the first term of \eqref{eq:three} is $1$, and on $S_3$ the second is $1$.

\head{Words of lengths four to seven.}For a word $w$ of length $n\ge4$, let $p(t)$ mark the lengths
of prefixes of $w$ in $\operatorname{supp}P$, and let $s(t)$ mark the lengths of suffixes in
$\operatorname{supp}S$, with coefficient one for each marked length. Cutting $w$ into a prefix and
a suffix, either adjacent or separated by one letter, gives
\[
  [w]F=[t^n]\,p(t)(t-1)s(t).
\]
A marked length three requires a marked length two on the same side. Thus each set of lengths is
an initial interval or has polynomial $1+t^2+et^3$ with $e\in\{0,1\}$. For an interval
$\{0,\dots,r\}$ its polynomial times $t-1$ equals $t^{r+1}-1$. The other factor has degree at most
three, so the coefficient at $n\ge4$ is nonnegative.

If both sides have a gap, denote their two bits by $e,f$. At degrees $4,5,6,7$ the coefficients
are, in order,
\[
  e+f-1,\qquad 1-e-f,\qquad e+f-ef,\qquad ef.
\]
For $w=xyzh$ of length four the gaps force $xy\in P_2$, $x\notin U$, $zh\in S_2$, $h\notin V$.
Since both pairs lie in $H$, $y\in V$ and $z\in U$. Hence $e=[yz\notin D]$ and $f=[yz\in D]$,
giving $e+f=1$. At length five, $e=1$ requires the middle letter in $U$ and $f=1$ requires it in
$V$, so $e+f\le1$. The last two coefficients are nonnegative for all bits. All degrees above seven
vanish. This proves the proposition.
\end{proof}

\section{Cardinality and total length}

For a word set $W$ let $\rev W$ consist of its words read backwards. Reversing a certificate
exchanges its two factors and reverses their words. Sending the other letters to zero restricts a
certificate to a subalphabet. Both operations preserve its nonnegativity and its constant
coefficients.

\head{At most three words.}The formalized three-word theorem \cite{ThreeWord} states that every
uniquely decodable set of exactly three words on a finite alphabet is contained, on that alphabet,
in a finite maximal code with a unique positive factorization as defined above. If $|C|<3$, adjoin
$3-|C|$ fresh letters as individual codewords. In a parsing the positions of these letters are
prescribed, and each intervening segment has at most one $C$-parsing. Thus the enlarged set is a
three-word code. Restrict a certificate of its positive completion to the original alphabet. This
proves the cardinality bound, including the empty code.

\begin{lemma}[Shortening by prefix quotients]\label{lem:shortening}
Let $H(W)$ be the smallest set containing a finite $W\subseteq A^+$ and closed under
$(u,uv)\mapsto v$ for $v\ne\eps$. If a letter occurring in a code $W$ is absent from $H(W)$, there
is a code $W'$ on a finite alphabet with $\tau(W')<\tau(W)$ whose positively factorizing finite
maximal completions give such completions of $W$.
\end{lemma}

\begin{proof}
All elements of $H(W)$ are nonempty suffixes of elements of $W$. Let $B$ consist of the elements
of $H(W)$ with no shorter prefix in $H(W)$. Then $B$ is prefix-free. To parse a member of $H(W)$
into $B$, remove its shortest prefix in $H(W)$. Every nonempty remainder again lies in $H(W)$, so
this process reaches the empty word. The parsing is unique because distinct elements of $B$
cannot start at the same cut. A word of $W$ containing the missing letter uses a block of length
at least two.

Take the finite tree whose internal vertices $I$ are the proper prefixes of $B$, and give every
internal vertex all its $A$-children. Its leaves $E$ contain $B$. Regard the leaves as distinct
letters of a new alphabet and let $\phi$ replace each such letter by its word on $A$. The leaves
are prefix-free, so $\phi$ is injective. The $B$-parsings of $W$ form a code $W'$ over $E$, since
distinct parsings of an equal concatenation over $W'$ would give distinct $W$-parsings of its
image. The block of length at least two makes the total length strictly smaller.

For clarity, $\chi_E$ here denotes the leaf polynomial on $A$. Counting the nonroot vertices of the
full tree by their parents gives $\chi_E-1=\chi_I(\chi_A-1)$. If a positive completion $M'$ of
$W'$ on the leaf alphabet has certificate factors $P',S'$, substituting $\phi$ gives
\[
  \chi_{\phi(M')}-1=\phi(P')\chi_I(\chi_A-1)\phi(S').
\]
This is a positive certificate on $A$, and its code contains $W$.
\end{proof}

\begin{lemma}[Finite bound for the two closures]\label{lem:closures}
Let $W$ be a code with $|W|\ge4$, $\lambda(W)\ge4$ and $\tau(W)\le10$. At least one of $H(W)$ and
$H(\rev W)$ omits a letter occurring in $W$.
\end{lemma}

\begin{proof}
Put $n=\tau(W)$, so $7\le n\le10$. Sort the words by length and concatenate them. Up to renaming,
after $k$ distinct letters have appeared, the next position contains one of those letters or the
next unused letter. Cutting all such strings at the partial sums of all nondecreasing length
profiles covers every $W$. Discard repeated words. Ties between lengths may represent the same set
more than once, which does not affect a zero count.

Let $T$ be the letters used. Since $H(W)$ consists of suffixes, the inclusion
$T\subseteq H(W)\cap H(\rev W)$ requires each letter to occur both initially and finally in $W$.
The columns successively impose this endpoint condition, the two closure conditions, and the
residual test of Lemma~\ref{lem:residual}.
\begin{center}
\begin{tabular}{rrrr}
$n$ & All letters at both ends & Both closures contain $T$ & Codes among these\\\hline
7 & 97 & 81 & 0\\
8 & 909 & 611 & 0\\
9 & 7\,106 & 3\,845 & 0\\
10 & 51\,910 & 24\,236 & 0
\end{tabular}
\end{center}
The complete calculation below implements $\partial$ as \texttt{tails}. Its fixed-point iterations
terminate because only nonempty suffixes of the words being tested are added. The two generators
enumerate the profiles and letter patterns just described. The last column proves the lemma.
\end{proof}

\head{Completion of the proof of Theorem~\ref{thm:main}.}The first two bounds follow from the
three-word result and Proposition~\ref{prop:short}. For the third, use strong induction on
$\tau(Y)\le10$ simultaneously over all finite alphabets. Outside the first two bounds,
Lemma~\ref{lem:closures} and, if necessary, reversal allow Lemma~\ref{lem:shortening} to shorten
$Y$. The induction hypothesis supplies the positive completion on the leaf alphabet, and the lemma
transports it back. Finally $X=Y_{2,2}$, so Theorem~\ref{thm:family} excludes every finite maximal
completion of $X$. Its word lengths are $1,2,4,4$, so the three lower bounds are attained on
every finite alphabet containing $a,b$.

\section*{Finite calculation for Lemma~\ref{lem:closures}}

\begin{Verbatim}[fontsize=\small]
from itertools import accumulate

def tails(U, V):
    return {v[len(u):] for u in U for v in V
            if u != v and v.startswith(u)}

def close(S, step):
    while True:
        T = S | step(S)
        if T == S:
            return S
        S = T

def hull(W):
    return close(W, lambda T: tails(T, T))

def code(W):
    def step(T):
        return tails(T, W) | tails(W, T)
    D = close(tails(W, W), step)
    return W.isdisjoint(D)

def profiles(n, lo=1):
    if n == 0:
        yield ()
    for k in range(lo, n + 1):
        for r in profiles(n - k, k):
            yield (k,) + r

def patterns(n, p='', k=0):
    if len(p) == n:
        yield p
        return
    for j in range(k + 1):
        q = p + chr(97 + j)
        yield from patterns(n, q, max(k,j+1))

def count(n):
    ends = both = codes = 0
    for L in profiles(n):
        if len(L) < 4 or L[-1] < 4:
            continue
        cuts = list(accumulate((0,) + L))
        spans = list(zip(cuts, cuts[1:]))
        for p in patterns(n):
            W = {p[i:j] for i,j in spans}
            A = set(p)
            if len(W) != len(L):
                continue
            if {w[0] for w in W} != A:
                continue
            if {w[-1] for w in W} != A:
                continue
            ends += 1
            R = {w[::-1] for w in W}
            if A <= hull(W) and A <= hull(R):
                both += 1
                codes += int(code(W))
    return ends, both, codes

for n in range(7, 11):
    print(n, *count(n))
\end{Verbatim}

\begin{thebibliography}{2}
\bibitem{Reutenauer} C.~Reutenauer. Noncommutative factorization of variable-length codes.
  \emph{Journal of Pure and Applied Algebra} 36 (1985), 167--186. Corollary of Section~2,
  p.~172. \href{https://doi.org/10.1016/0022-4049(85)90069-6}{doi:10.1016/0022-4049(85)90069-6}.
\bibitem{ThreeWord} Formal verification of the theory of codes. Lean declaration
  \texttt{ThreeWordCodes.three\_word\_positive\_completion}, in
  \texttt{ThreeWordCodes/MainTheorem.lean}, lines 16--19. The statement and the parsing predicate
  are in \texttt{Statement.lean} and \texttt{Definitions.lean} in the same directory. Repository
  revision \texttt{ff2b0011245e3eec688ea3ffe1181179958cf533}.
\end{thebibliography}

\end{document}
